% Modelo do livro XRiscs (XeLaTeX + pandoc). Formato 7 x 10 pol, tipografia clássica de LaTeX.
\documentclass[10pt,twoside,openright]{book}
\usepackage[paperwidth=7in,paperheight=10in,inner=0.95in,outer=0.8in,top=0.9in,bottom=0.95in,headsep=0.28in,footskip=0.42in]{geometry}
\usepackage{fontspec}
\usepackage{amsmath,amssymb,amsthm,mathtools}
\usepackage[warnings-off={mathtools-colon,mathtools-overbracket}]{unicode-math}
\setmainfont{Latin Modern Roman}[SmallCapsFont={Latin Modern Roman Caps}]
\setsansfont{Latin Modern Sans}
\setmonofont{DejaVu Sans Mono}[Scale=0.78]
\setmathfont{Latin Modern Math}
\usepackage{polyglossia}
\setdefaultlanguage[variant=brazilian]{portuguese}
\usepackage{microtype}
\usepackage{xcolor}
\definecolor{tinta}{HTML}{1B2E4B}
\definecolor{acento}{HTML}{2A78D6}
\definecolor{laranja}{HTML}{C9562B}
\definecolor{fundocod}{HTML}{F5F6F8}
\definecolor{fundosai}{HTML}{FCFCFA}
\definecolor{cinza}{HTML}{7A7A74}
\usepackage{graphicx}
\usepackage{adjustbox}
\usepackage{booktabs,longtable,array,calc}
\usepackage{tikz}
\usepackage{titlesec}
\usepackage{fancyhdr}
\usepackage{titletoc}
\usepackage{enumitem}
\setlist{itemsep=2pt,topsep=4pt}
\usepackage{fvextra}
\usepackage[most]{tcolorbox}
\usepackage{etoolbox}
\usepackage[round,authoryear]{natbib}
\setlength{\bibsep}{3pt plus 1pt}
\usepackage[hidelinks,unicode,bookmarksnumbered]{hyperref}
\usepackage{bookmark}
\setcounter{secnumdepth}{0}
\setcounter{tocdepth}{0}
\providecommand{\tightlist}{\setlength{\itemsep}{0pt}\setlength{\parskip}{0pt}}
\providecommand{\pandocbounded}[1]{#1}
\setlength{\emergencystretch}{2.5em}
\numberwithin{equation}{chapter}

% ---------- teoremas ----------
\newtheoremstyle{livro}{8pt}{8pt}{\itshape}{}{\bfseries}{.}{0.5em}{\thmname{#1}\thmnumber{ #2}\thmnote{ \normalfont(#3)}}
\newtheoremstyle{livrodef}{8pt}{8pt}{\normalfont}{}{\bfseries}{.}{0.5em}{\thmname{#1}\thmnumber{ #2}\thmnote{ \normalfont(#3)}}
\theoremstyle{livro}
\newtheorem{teorema}{Teorema}[chapter]
\newtheorem{proposicao}[teorema]{Proposição}
\newtheorem{lema}[teorema]{Lema}
\newtheorem{corolario}[teorema]{Corolário}
\theoremstyle{livrodef}
\newtheorem{definicao}[teorema]{Definição}
\newtheorem{exemplo}[teorema]{Exemplo}
\newtheorem{observacao}[teorema]{Observação}
\renewcommand{\proofname}{Demonstração}
\renewcommand{\qedsymbol}{\(\blacksquare\)}

% ---------- capítulos e seções ----------
\titleformat{\chapter}[display]
  {\normalfont\filleft}
  {{\large\scshape\color{cinza}\chaptertitlename}\hspace{0.4em}{\fontsize{64}{64}\selectfont\color{tinta}\thechapter}}
  {14pt}
  {\fontsize{24}{29}\selectfont\bfseries\color{tinta}\filleft}
  [\vspace{6pt}{\color{tinta}\titlerule[0.8pt]}]
\titlespacing*{\chapter}{0pt}{0.3in}{32pt}
\titleformat{name=\chapter,numberless}[display]{\normalfont\filleft}{}{0pt}
  {\fontsize{24}{29}\selectfont\bfseries\color{tinta}\filleft}[\vspace{6pt}{\color{tinta}\titlerule[0.8pt]}]
\titlespacing*{name=\chapter,numberless}{0pt}{0.3in}{32pt}
\titleformat{\section}{\normalfont\Large\bfseries\color{tinta}\raggedright}{}{0pt}{}
\titleformat{\subsection}{\normalfont\large\bfseries\color{tinta}\raggedright}{}{0pt}{}
\titleformat{\subsubsection}{\normalfont\normalsize\bfseries\itshape\raggedright}{}{0pt}{}
\titlespacing*{\section}{0pt}{18pt plus 4pt}{7pt}
\titlespacing*{\subsection}{0pt}{13pt plus 3pt}{5pt}
\titleformat{\part}[display]{\thispagestyle{empty}\centering\normalfont\color{tinta}}
  {\Large\scshape Parte\ \thepart}{16pt}{\fontsize{28}{34}\selectfont\bfseries}[\vspace{10pt}{\color{cinza}\rule{0.3\linewidth}{0.5pt}}]
\renewcommand{\thepart}{\Roman{part}}
\titlecontents{part}[0pt]{\addvspace{14pt}\normalfont\bfseries\color{tinta}}{Parte \thecontentslabel\quad}{}{\hfill\contentspage}[\addvspace{3pt}]
\titlecontents{chapter}[2.2em]{\addvspace{3pt}\normalfont}{\contentslabel{2.2em}}{}{\titlerule*[0.7pc]{.}\contentspage}

% ---------- cabeçalho e rodapé ----------
\pagestyle{fancy}
\fancyhf{}
\renewcommand{\headrulewidth}{0pt}
\renewcommand{\chaptermark}[1]{\markboth{#1}{}}
\renewcommand{\sectionmark}[1]{\markright{#1}}
\fancyhead[LE]{\small\thepage\hspace{1.2em}\scshape XRiscs --- Derivativos com Python}
\fancyhead[RO]{\small\scshape\nouppercase{\leftmark}\hspace{1.2em}\upshape\thepage}
\fancypagestyle{plain}{\fancyhf{}\fancyfoot[C]{\small\thepage}}

% ---------- código e saídas ----------
$if(highlighting-macros)$
$highlighting-macros$
$endif$
\RecustomVerbatimEnvironment{Highlighting}{Verbatim}{commandchars=\\\{\},breaklines,breakanywhere,fontsize=\footnotesize}
\renewenvironment{Shaded}{\begin{tcolorbox}[enhanced,breakable,colback=fundocod,colframe=fundocod,
  borderline west={1.2pt}{0pt}{tinta!60},arc=0pt,outer arc=0pt,boxrule=0pt,left=6pt,right=4pt,top=2pt,bottom=2pt,before skip=6pt,after skip=6pt]}{\end{tcolorbox}}
\DefineVerbatimEnvironment{Saida}{Verbatim}{breaklines,breakanywhere,fontsize=\scriptsize,formatcom=\color{black!80}}
\newenvironment{saida}{\begin{tcolorbox}[enhanced,breakable,colback=fundosai,colframe=black!15,boxrule=0.4pt,arc=0pt,
  left=6pt,right=4pt,top=1pt,bottom=1pt,before skip=2pt,after skip=8pt]}{\end{tcolorbox}}
\newtcolorbox{solucao}{enhanced,breakable,colback=white,colframe=tinta!70,boxrule=0.5pt,arc=0pt,
  title={\bfseries Solução},coltitle=white,colbacktitle=tinta!70,fonttitle=\small,left=6pt,right=6pt}
\newtcolorbox{nota}{enhanced,breakable,colback=acento!5,colframe=acento!70,boxrule=0.5pt,arc=0pt,left=6pt,right=6pt}
\newcommand{\figura}[2]{\par\smallskip{\centering\includegraphics[width=#2\linewidth]{#1}\par}\smallskip}
\newcommand{\tabelalivro}[1]{\par\smallskip{\centering\small\begin{adjustbox}{max width=\linewidth}#1\end{adjustbox}\par}\smallskip}
% abertura da seção teórica de cada capítulo
\newcommand{\aberturateoria}{\par\noindent{\color{cinza}\rule{\linewidth}{0.4pt}}\par\nobreak}

\begin{document}
\frontmatter
\pagestyle{empty}
% ---------- capa ----------
\begin{tikzpicture}[remember picture,overlay]
  \node[anchor=center,inner sep=0] at (current page.center) {\includegraphics[width=\paperwidth,height=\paperheight]{capa_fundo.pdf}};
  \node[anchor=north west,text width=5.6in,align=left,text=white] at ([xshift=0.75in,yshift=-1.05in]current page.north west)
    {{\fontsize{50}{54}\selectfont\bfseries XRiscs}\\[16pt]
     {\fontsize{26}{30}\selectfont Derivativos com Python}\\[14pt]
     {\color{white!80}\rule{1.6in}{0.6pt}}\\[12pt]
     {\Large\itshape Precificação, calibração e hedge em Python,\\ do Black-Scholes ao smile, com auxílio de IA}};
  \node[anchor=south west,text=white] at ([xshift=0.75in,yshift=1.15in]current page.south west) {\fontsize{20}{24}\selectfont\scshape Walter C. Neto};
  \node[anchor=south west,text=white!75] at ([xshift=0.75in,yshift=0.8in]current page.south west) {\small Edição $edicao$ \,·\, $data$};
\end{tikzpicture}
\cleardoublepage
% ---------- folha de rosto ----------
\begin{center}
\vspace*{1.3in}{\fontsize{34}{40}\selectfont\bfseries\color{tinta} XRiscs}\\[12pt]
{\LARGE Derivativos com Python}\\[10pt]
{\large\itshape Precificação, calibração e hedge em Python,\\ do Black-Scholes ao smile, com auxílio de IA}\\[1.3in]
{\Large\scshape Walter C. Neto}\\[2.1in]
{\small waltercneto.github.io/XRiscs}
\end{center}
\newpage
% ---------- página de créditos ----------
\vspace*{\fill}
{\small\noindent
\textbf{XRiscs --- Derivativos com Python}\\
Walter C. Neto\\
Edição $edicao$, $data$. Gerada a partir dos notebooks do repositório
\url{https://github.com/walterCNeto/XRiscs} (commit $commit$).\par\medskip
\noindent Textos e figuras: Creative Commons Atribuição 4.0 Internacional (CC BY 4.0).
Código: licença MIT. Você pode copiar, adaptar e redistribuir, desde que cite a fonte.\par\medskip
\noindent Textos, código e laboratórios desenvolvidos com auxílio de IA (Claude, da Anthropic).\par\medskip
\noindent Este é um material de estudo independente. A sequência dos Capítulos 1 a 16 acompanha
\citet{hilpisch2015}; não há afiliação com o autor nem com a editora, e nenhum trecho daquele livro é reproduzido.
Todos os dados são sintéticos, exceto as taxas dos T-bills do Capítulo 17 (U.S. Treasury, domínio público).\par\medskip
\noindent Cada capítulo corresponde a um notebook executável. As verificações automáticas \texttt{ok(...)}
que aparecem nas saídas foram todas aprovadas na geração desta edição ($verificacoes$ no total).\par\medskip
\noindent Capa: trajetórias simuladas de um movimento browniano geométrico, com as bandas de \(\pm1\) e \(\pm2\) desvios-padrão do log-preço.\par}
\cleardoublepage
\pagestyle{fancy}
\tableofcontents
\cleardoublepage
$body$
\backmatter
\cleardoublepage
\phantomsection
\addcontentsline{toc}{chapter}{Referências}
\renewcommand{\bibname}{Referências}
\nocite{*}
\bibliographystyle{plainnat}
\bibliography{referencias}
\end{document}
